\begin{helper}
In the fixed history cell, assume \(u_B\le u_R\).  With
\[
a=f(\ell_R,\ell_B)-10\varepsilon_1 k^2,
\]
there are natural numbers \(N_R,N_B\) such that
\[
N_R\le 3\varepsilon_1 k^2,\qquad N_B\le k^2,
\]
and the conditional probability of the fixed-count target is at most
\[
\binom{N_R}{u_R}p^{u_R}\binom{N_B}{u_B}p^{u_B}
(1-p)^{\max(0,a-u_R-u_B)}.
\]
Here red is the later-color support controlled by the exceptional sets, and
blue is the first-color support controlled by the \(k\)-set projection
universe.
\end{helper}
\begin{proof}
Let \(H\) be the history cell and let \(T_{u_R,u_B}\) be the target event:
regularity holds, the final \(I\)-set is independent, and the red and blue
present staged-open supports have sizes \(u_R\) and \(u_B\).  Put
\[
L=\max(0,a-u_R-u_B).
\]
All terminal randomness is contained in the finite terminal coordinate set,
which we denote by \(U\), and the complete terminal product law is available through
\(\noderef{FixedSetHistoryCellTerminalProductPartition}\).

First use \(\noderef{FixedSetHistoryCellGlobalRedExceptionalWitness}\).  It
turns the branch-dependent red exceptional sets into one global number
\(N_R\le3\varepsilon_1k^2\).  For each complete blue terminal branch, every
regular independent completion has its red present staged-open support inside
one exceptional set of cardinality at most \(N_R\).  For each complete blue
branch we pad its exceptional set to \(N_R\) slots and label a red support by
a \(u_R\)-subset of those slots.  Thus the red support has at most
\(\binom{N_R}{u_R}\) labels, and a label is decoded to an actual red terminal
support only after the complete blue trace is fixed.

Next let \(U_B\) be the finite universe of blue projected terminal pairs
coming from pairs of vertices in \(I\).  Since \(|I|=k\), this universe has
cardinality at most \(k^2\); call its cardinality \(N_B\).  The blue present
support of every sample in \(H\cap T_{u_R,u_B}\) is a \(u_B\)-subset of this
universe, so there are at most \(\binom{N_B}{u_B}\) blue support labels.

Fix a blue support label \(B\) and a red label \(J\).  The branch label
\(\beta\) records a complete blue terminal trace
\(\operatorname{blueTrace}(\beta)\), a subset of the blue terminal part of
\(U\): for every blue terminal coordinate \(e\), a sample in the branch
\(\beta\) has \(X_e=1\) exactly when
\(e\in\operatorname{blueTrace}(\beta)\).  The cover and disjointness of the
complete blue traces are supplied by
\(\noderef{FixedSetHistoryCellAdaptiveBranchPartition}\), and
\(\noderef{FixedSetHistoryCellRedBranchTranscriptCells}\) chooses a branch
label uniquely from this terminal blue truth table: two labels that agree on
all blue coordinates of \(U\) are equal, and
\(\operatorname{blueTrace}(\beta)\) contains no red or nonterminal
coordinates.  Once \(\beta\) is fixed, the padded exceptional-set ordering
decodes \(J\) to a red terminal support \(R_{\beta,J}\).  If \(J\) uses a
padding slot, or if the decoded set does not have cardinality \(u_R\), the
corresponding cells are declared empty and receive weight \(0\).  Otherwise
\(R_{\beta,J}\) is a \(u_R\)-subset of the exceptional set for the
representative branch of \(\beta\).  The global red witness then says that
every regular independent completion in this blue trace has its red present
staged-open support equal to one of the decoded \(R_{\beta,J}\)'s.

For a fixed nonempty pair \((B,J)\), refine each blue trace by the canonical
prefix scan.  The scan reads the terminal order from left to right.  At the
moment a coordinate is inspected, the prefix-safety hypothesis in the active
statement says that staged-openness, touchedness, same-color closedness, and
recordedness are functions only of the history and the earlier terminal
values.  Therefore the scan can be encoded by the finite transcript
\[
\tau=(P_\tau,A_\tau,Z_\tau),
\]
where \(P_\tau\) is the set of terminal coordinates whose value is fixed
present by the complete trace, the decoded support, or the scan; \(A_\tau\)
is the analogous absent set; and \(Z_\tau\subseteq A_\tau\) is the selected
absence set produced by the canonical scan.  The set \(Z_\tau\) is not
assumed to be red-only.  A selected absence may be a red coordinate, or it may
be an absent blue coordinate already fixed by the complete blue trace.  These
are coordinate sets, not ordered role lists: if a
terminal coordinate occurs both as a support coordinate and later as a scan
datum, it is inserted once into \(P_\tau\), and similarly for \(A_\tau\).
The transcript is admitted only when
\[
B\cup R_{\beta,J}\subseteq P_\tau,\qquad
Z_\tau\subseteq A_\tau,\qquad
P_\tau\cap A_\tau=\emptyset .
\]
The canonical absence lemma
\(\noderef{FixedSetHistoryCellCanonicalAbsenceSelection}\) gives
\(|Z_\tau|\ge L\), disjointness from the present red and blue supports, and
the prefix measurability of the selected coordinates in the residual scan.
The later summation will extract the factors for all coordinates of
\(Z_\tau\) from the combined fixed cylinder at once; if a member of
\(Z_\tau\) is blue, the factor is the same \((1-p)\)-factor already present
in the complete blue trace cylinder, not a second charge.

The resulting cells \(G_{B,J,\beta,\tau}\) cover \(H\cap T_{u_R,u_B}\).
Indeed a target sample supplies \(B\) as its blue present staged-open support
and \(\beta\) as its complete blue terminal trace; the red exceptional
witness supplies the label \(J\) whose decoded support is its red present
staged-open support; and the deterministic prefix scan supplies \(\tau\).
Conversely, every nonempty cell forces all red and blue present staged-open
coordinates to be in \(R_{\beta,J}\cup B\), requires the selected coordinates
of \(Z_\tau\) to be absent, and keeps all fixed terminal values inside the
terminal universe.  If a sample belonged to two cells with the same \(B,J\),
the complete-trace disjointness first forces the same blue trace, and the
terminal-blue uniqueness clause in
\(\noderef{FixedSetHistoryCellRedBranchTranscriptCells}\) forces the same
\(\beta\): agreement on every blue coordinate of \(U\), not equality of an
unrestricted auxiliary set, is the uniqueness criterion.  With \(\beta\)
fixed, prefix-safety makes the scan deterministic from the earlier terminal
values, hence the same sample gives the same transcript \(\tau\).  Thus
there is no multiplicity inside the fixed support-label summation.

For each nonempty cell set its cylinder weight to
\[
w_{B,J,\beta,\tau}=
p^{|P_\tau|}(1-p)^{|A_\tau|},
\]
and set the weight to \(0\) for empty cells.  Since \(P_\tau\) and
\(A_\tau\) are disjoint subsets of the terminal set, the terminal product law
and \(\noderef{FixedSetHistoryCellFixedCylinderCharge}\) give
\[
\Pr(T_{u_R,u_B}\cap H\cap G_{B,J,\beta,\tau})
\le
\Pr(H)\,w_{B,J,\beta,\tau}.
\]
The use of coordinate sets here is essential: repeated appearances of one
terminal coordinate in different logical roles do not multiply the Bernoulli
factor, because the cylinder charges one coordinate value exactly once.

It remains to verify, for this fixed pair \((B,J)\), the fixed-support
branch/transcript weight sum required by
\(\noderef{FixedSetHistoryCellRedFixedCountCellSummation}\).  The
verification is deliberately carried out on realized cells first and only
then extended to the formal index set.  Let
\(\mathcal C_{B,J}\) be the finite set of realized, nonempty
trace/transcript cells.  Formal pairs outside \(\mathcal C_{B,J}\) include
all padding labels, bad decodes of the wrong size or color, prefix
transcripts inconsistent with the complete assignment, and empty cylinders.
Every non-realized formal pair is assigned weight \(0\); no geometric
assertion is required for it.
For \(c=(\beta,\tau)\), write
\[
R_c=R_{\beta,J},\qquad
Z_c=Z_\tau,\qquad
P_c=P_\tau,\qquad
A_c=A_\tau .
\]
The cell construction in
\(\noderef{FixedSetHistoryCellRedBranchTranscriptCells}\) gives, in every
nonempty cell,
\[
|R_c|=u_R,\qquad |B|=u_B,\qquad
B\cup R_c\subseteq P_c,\qquad Z_c\subseteq A_c,
\]
\[
P_c\cap A_c=\emptyset,\qquad |Z_c|\ge L.
\]
It also records the complete terminal blue trace in the transcript: for
every blue terminal coordinate \(e\), \(e\) belongs to the branch trace
exactly when \(e\in P_c\), and otherwise \(e\in A_c\).  In particular the
fixed blue support \(B\) is part of the branch trace and is charged inside
\(P_c\).
The weights supplied by
\(\noderef{FixedSetHistoryCellRedCellCylinderMass}\) satisfy
\[
0\le w_c\le p^{|P_c|}(1-p)^{|A_c|}.
\]
These are the only geometric facts used in the fixed-support summation below.
In particular, \(Z_c\) is only an absent selected set of size at least \(L\);
we do not assume that it consists of red coordinates.

Put \(q=1-p\).  Then \(1/2\le q\le1\).  For a nonempty cell
\(c=(\beta,\tau)\), let
\[
M_c=B\cup R_c.
\]
The sets \(B\) and \(R_c\) are disjoint because \(B\) is blue-only and
\(R_c\) is red-only, so \(|M_c|=u_B+u_R\).  Also
\[
M_c\subseteq P_c,\qquad Z_c\subseteq A_c,\qquad P_c\cap A_c=\emptyset.
\]
Thus the mandatory present coordinates \(M_c\) and the selected absent
coordinates \(Z_c\) are pairwise disjoint pieces of the same terminal
cylinder.  The remaining factors in
\(P_c\setminus M_c\) and \(A_c\setminus Z_c\) are precisely the auxiliary
values used to choose the complete blue trace and to run the prefix scan.
Those auxiliary values will be summed as residual product-tree factors; they
are not counted by a separate multiplicity bound.

At this point the proof needs fixed-\((B,J)\) released residual blocks, not
merely the forward compatibility and functionality statements exported by
\(\noderef{FixedSetHistoryCellRedBranchTranscriptCells}\).  The
branch/transcript construction also supplies the fixed residual
compatibility inputs: residual-cylinder completeness after deleting
selected absences and reinserting \(B\cup R_{\beta,J}\), and single-owner
uniqueness for one residual assignment.  Apply
\(\noderef{FixedSetHistoryCellFixedBJResidualScanInterface}\) to these
inputs to construct one residual scan and one reconstructed full scan tied
to the fixed labels.  Then apply
\(\noderef{FixedSetHistoryCellFixedBJResidualScanData}\), which extracts
from that interface exactly the two facts needed here:
residual-cylinder completeness and single-owner uniqueness for the same
residual assignment, even when the two candidate reconstructions use
different selected sets or fixed supports.  These are not consequences of
forward compatibility alone.  With those scan facts in place, apply
\(\noderef{FixedSetHistoryCellFixedBJResidualReleasedBlockData}\).
It supplies blocks
\(\mathcal R_{\beta,\tau}\) whose membership is exactly reconstructed
compatibility for
\[
(A_{\rm res}\setminus Z_\tau)\cup B\cup R_{\beta,J},
\]
whose blocks are disjoint for distinct realized pairs, and whose blocks
contain the residual cylinder obtained by releasing
\(B\cup R_{\beta,J}\) and \(Z_\tau\).

Now apply
\(\noderef{FixedSetHistoryCellFixedBJResidualProductTreeOwnerPackage}\) to
these released blocks together with the realized geometry just listed:
\(B\) has size \(u_B\) and is contained in the terminal set, each realized
\(R_{\beta,J}\) has size \(u_R\) and is red-only, \(B\cup R_{\beta,J}
\subseteq P_\tau\), \(Z_\tau\subseteq A_\tau\), \(P_\tau\cap A_\tau=\emptyset\),
and the complete blue trace is encoded by \(P_\tau,A_\tau\) on all blue
terminal coordinates.  For the fixed labels \(B,J\) the package supplies one
owner function
\[
\operatorname{owner}_{B,J}:2^U\to
  \{\hbox{branch/transcript pairs}\}\cup\{\bot\},
\]
not a candidate-dependent reconstruction rule.  For every realized
\((\beta,\tau)\), its owner fiber contains the residual cylinder
\[
\{A_{\rm res}\subseteq U:
P_\tau\setminus(B\cup R_{\beta,J})\subseteq A_{\rm res},\ 
A_{\rm res}\cap(A_\tau\setminus Z_\tau)=\emptyset\}.
\]
The owner uniqueness is justified by the released-block disjointness: if
one assignment lay in the blocks for two realized pairs, the disjointness
invariant would force equality of both the branch label and the transcript.
The selected coordinates \(Z_\tau\) are released from this residual tree
uniformly; they may be red or blue, and their \((1-p)\)-factors are charged
only once.

The same package also returns the normalized residual leaf certificate
needed by \(\noderef{FixedSetHistoryCellMixedBranchTranscriptWeightPackage}\).
For the fixed \(B,J\) it gives a finite leaf type, nonnegative leaf masses
with total mass at most \(1\), injectivity of the leaf map on realized
branch/transcript cells, and for every realized \((\beta,\tau)\) the
domination
\[
p^{|P_\tau|-u_R-u_B}(1-p)^{|A_\tau|-|Z_\tau|}
  \le \nu(\iota(\beta,\tau)).
\]
This is exactly the residual subprobability estimate for the auxiliary
values in \(P_\tau\setminus(B\cup R_{\beta,J})\) and
\(A_\tau\setminus Z_\tau\).

For a realized cell \(c=(\beta,\tau)\), set
\[
R_c=R_{\beta,J}.
\]
The coordinate sets charged as mandatory factors are disjoint:
\[
B\cap R_c=\emptyset,\qquad
Z_\tau\cap(B\cup R_c)=\emptyset .
\]
The first disjointness is by color, and the second follows from
\[
B\cup R_c\subseteq P_\tau,\qquad Z_\tau\subseteq A_\tau,\qquad
P_\tau\cap A_\tau=\emptyset .
\]
Thus each realized cell cylinder has a unique factorization
\[
p^{|P_\tau|}q^{|A_\tau|}
=
p^{u_R}p^{u_B}q^{|Z_\tau|}\,\rho_{\beta,\tau},
\]
where \(\rho_{\beta,\tau}\) is the product of all remaining terminal factors
from \(P_\tau\setminus(B\cup R_c)\) and \(A_\tau\setminus Z_\tau\).  If a
selected absent coordinate is blue, its factor is the same \(q\)-factor
already present in the complete blue trace; it is removed once into
\(q^{|Z_\tau|}\), and no second factor is introduced.  If an absent blue
coordinate is not selected, its \(q\)-factor stays in the residual factor.
The fixed-\((B,J)\) leaf certificate just cited says that these residual
factors can be charged to injective leaves of a normalized product tree, so
their total contribution over realized cells is at most \(1\).

The cell-mass helper
\(\noderef{FixedSetHistoryCellRedCellCylinderMass}\) gives
\[
0\le w_{B,J,\beta,\tau}
\le p^{|P_\tau|}q^{|A_\tau|}
\]
for realized cells, and every non-realized formal pair has weight \(0\).
Since \(|Z_\tau|\ge L\) by
\(\noderef{FixedSetHistoryCellCanonicalAbsenceSelection}\) and \(q=1-p\le1\),
we have \(q^{|Z_\tau|}\le q^L\).  Therefore
\[
\begin{aligned}
\sum_{\beta,\tau}w_{B,J,\beta,\tau}
&=
\sum_{\beta,\tau:\mathsf{Realized}(\beta,\tau)}
  w_{B,J,\beta,\tau} \\
&\le
\sum_{\beta,\tau:\mathsf{Realized}(\beta,\tau)}
  p^{u_R}p^{u_B}q^{|Z_\tau|}\rho_{\beta,\tau} \\
&\le
p^{u_R}p^{u_B}q^L
\sum_{\beta,\tau:\mathsf{Realized}(\beta,\tau)}
  \rho_{\beta,\tau} \\
&\le p^{u_R}p^{u_B}q^L .
\end{aligned}
\]
This is exactly the finite branch/transcript overcount estimate required for
the fixed support labels.  It uses the complete blue trace as part of the
same terminal product tree, rather than conditioning on a blue trace and then
trying to prove that the selected absences are red-only.

Applying
\(\noderef{FixedSetHistoryCellMixedBranchTranscriptWeightPackage}\) with the
cell geometry, non-realized-cell exclusion, compatibility disjointness, and
the residual leaf certificate just constructed supplies cell weights with the
same pointwise cylinder bound and fixed-\((B,J)\) sum bound.  The hypotheses of
\(\noderef{FixedSetHistoryCellRedFixedCountCellSummation}\) are therefore
satisfied.  Its blue support-label family is the \(u_B\)-subsets of the
\(I\)-projection universe, so it has size at most \(\binom{N_B}{u_B}\).  Its
red support-label family is the padded family of \(u_R\)-slot choices from
the global red witness, so it has size at most \(\binom{N_R}{u_R}\).  The
cell cover, the cylinder mass estimate, and the fixed-support
branch/transcript sum are exactly the conclusions established above.
Applying that helper yields
\[
\Pr(T_{u_R,u_B}\mid H)\le
\binom{N_R}{u_R}p^{u_R}\binom{N_B}{u_B}p^{u_B}(1-p)^L,
\]
which is the displayed bound.  Its proof also handles the zero-history-mass
case through \(\noderef{TwoBiteConditionalProbability}\).
\end{proof}
